\documentclass[tikz,border=8pt]{standalone}
\usepackage{tikz}
\usetikzlibrary{arrows.meta, positioning, calc}
\usepackage{amsmath}
\usepackage{pgfplots}
\pgfplotsset{compat=1.18}
\usepackage{ctex}
\begin{document}
\begin{tikzpicture}

% --- Left panel: X_n ->_d X (PDF evolution) ---
\begin{axis}[
  name=ax1,
  width=6cm, height=5.5cm,
  xmin=-5, xmax=5,
  ymin=0, ymax=0.50,
  xlabel={$x$},
  ylabel={$f(x)$},
  title={$X_n \xrightarrow{d} X$},
  axis lines=left,
  xtick={-4,-2,0,2,4},
  xticklabel style={font=\scriptsize},
  yticklabel style={font=\scriptsize},
  ytick={0,0.1,0.2,0.3,0.4},
  grid=major, grid style={gray!30},
  clip=true,
  legend pos=north east,
  legend style={font=\tiny}
]
  % n=1: wide spread
  \addplot[blue!30, thin, domain=-5:5, samples=100]
    {exp(-x^2/(2*4))/sqrt(2*3.14159*4)};
  \addlegendentry{$n=1$}
  % n=5: medium
  \addplot[blue!55, thin, domain=-5:5, samples=100]
    {exp(-x^2/(2*2))/sqrt(2*3.14159*2)};
  \addlegendentry{$n=5$}
  % n=20: narrower
  \addplot[blue!80, thick, domain=-5:5, samples=100]
    {exp(-x^2/(2*1.2))/sqrt(2*3.14159*1.2)};
  \addlegendentry{$n=20$}
  % Limit X ~ N(0,1)
  \addplot[red!70!black, ultra thick, dashed, domain=-5:5, samples=100]
    {exp(-x^2/2)/sqrt(2*3.14159)};
  \addlegendentry{$X\sim N(0,1)$}
\end{axis}

% Arrow 1: X_n ->_d X contributes to result
\draw[-Stealth, very thick, gray!60]
  ([xshift=3pt]ax1.east) -- ++(1.0cm, 0)
  node[midway, above, font=\scriptsize] {$\xrightarrow{d}$};

% --- Middle panel: Y_n ->_p c (scatter cloud narrowing to a point) ---
\begin{axis}[
  name=ax2,
  at={(ax1.east)}, anchor=west, xshift=1.4cm,
  width=5.5cm, height=5.5cm,
  xmin=0, xmax=11,
  ymin=0.5, ymax=3.5,
  xlabel={$n$},
  ylabel={$Y_n$},
  title={$Y_n \xrightarrow{P} c=2$},
  axis lines=left,
  xtick={1,3,5,8,10},
  xticklabel style={font=\scriptsize},
  yticklabel style={font=\scriptsize},
  ytick={1,1.5,2,2.5,3},
  grid=major, grid style={gray!30},
  clip=true
]
  % Scatter points narrowing toward c=2
  \addplot[only marks, mark=*, mark size=1.2pt, orange!80!black, opacity=0.8] coordinates {
    (1,1.0) (1,2.8) (1,1.5) (1,3.0) (1,1.2) (1,2.6)
    (2,1.3) (2,2.6) (2,1.6) (2,2.5) (2,1.4)
    (3,1.5) (3,2.4) (3,1.7) (3,2.3) (3,1.6)
    (5,1.7) (5,2.2) (5,1.9) (5,2.1) (5,1.8)
    (8,1.85) (8,2.12) (8,1.92) (8,2.08)
    (10,1.93) (10,2.06) (10,1.97) (10,2.03)
  };
  % Convergence bands
  \addplot[red!60, dashed, domain=1:10, samples=2] {2};
  \addplot[gray!60, dashed, thin, domain=1:10, samples=50] {2 + 1/x};
  \addplot[gray!60, dashed, thin, domain=1:10, samples=50] {2 - 1/x};
  \node[font=\scriptsize, red!70!black] at (axis cs:5,2.25) {$c=2$};
\end{axis}

% Arrow 2: result
\draw[-Stealth, very thick, gray!60]
  ([xshift=3pt]ax2.east) -- ++(1.0cm, 0)
  node[midway, above=4pt, font=\scriptsize] {Slutsky};

% --- Right panel: X_n + Y_n ->_d X + c (shifted distribution) ---
\begin{axis}[
  name=ax3,
  at={(ax2.east)}, anchor=west, xshift=1.6cm,
  width=6cm, height=5.5cm,
  xmin=-3, xmax=7,
  ymin=0, ymax=0.50,
  xlabel={$x$},
  ylabel={$f(x)$},
  title={$X_n+Y_n \xrightarrow{d} X+c$},
  axis lines=left,
  xtick={-2,0,2,4,6},
  xticklabel style={font=\scriptsize},
  yticklabel style={font=\scriptsize},
  ytick={0,0.1,0.2,0.3,0.4},
  grid=major, grid style={gray!30},
  clip=true,
  legend pos=north west,
  legend style={font=\tiny}
]
  % Original X ~ N(0,1) (faint)
  \addplot[blue!40, thin, dashed, domain=-3:7, samples=100]
    {exp(-x^2/2)/sqrt(2*3.14159)};
  \addlegendentry{$X\sim N(0,1)$}
  % Shifted: X + c ~ N(c, 1) = N(2, 1)
  \addplot[red!70!black, ultra thick, domain=-3:7, samples=100]
    {exp(-(x-2)^2/2)/sqrt(2*3.14159)};
  \addlegendentry{$X{+}c\sim N(2,1)$}
  % Shift arrow annotation
  \draw[-Stealth, thick, orange!80!black]
    (axis cs:0,0.30) -- (axis cs:2,0.30)
    node[midway, above, font=\scriptsize] {平移 $c$};
\end{axis}

% Slutsky theorem title at top
\node[font=\normalsize\bfseries] at ($(ax1.north)!0.5!(ax3.north)+(0,1.2cm)$)
  {Slutsky 定理：$X_n\xrightarrow{d}X,\ Y_n\xrightarrow{P}c\ \Rightarrow\ X_n{+}Y_n\xrightarrow{d}X{+}c$};

\end{tikzpicture}
\end{document}
